\documentclass[a4,11pt]{aleph-notas}
% Para obtener solo el pdf, compilar con pdfLaTeX.
% Para obtener el xml compilar con XeLaTeX.

% -- Paquetes adicionales
\usepackage[handout]{aleph-moodle}
% \usepackage{aleph-moodle}
\moodleregisternewcommands
% Todos los comandos nuevos deben ir luego del comando anterior
\usepackage{aleph-comandos}


% -- Datos
\institucion{Facultad de Ciencias Exactas, Naturales y Ambientales}
\carrera{Catálogo STEM}
\asignatura{Matemática III}
\tema{Cuestionario No. 1: Funciones vectoriales y geometría de curvas}
\autor{Andrés Merino}
\fecha{Periodo 2026-1}

\logouno[0.16\textwidth]{Logos/logoPUCE_04_ac}
\definecolor{colortext}{HTML}{1957d1}
\definecolor{colordef}{HTML}{1957d1}
\fuente{montserrat}


\begin{document}

\encabezado

\vspace*{-8mm}
\tableofcontents

%%%%%%%%%%%%%%%%%%%%%%%%%%%%%%%%%%%%%%%%
\section{Indicaciones}
%%%%%%%%%%%%%%%%%%%%%%%%%%%%%%%%%%%%%%%%

Se plantean bancos de preguntas orientados a evaluar el {criterio}: «Asocia las funciones de varias variables con su ley de asignación y su representación gráfica», correspondiente al {resultado de aprendizaje}: Identificar los conceptos, teoremas y propiedades de funciones vectoriales, derivadas parciales e integral vectorial.

%%%%%%%%%%%%%%%%%%%%%%%%%%%%%%%%%%%%%%%%
\section{Banco de preguntas}
%%%%%%%%%%%%%%%%%%%%%%%%%%%%%%%%%%%%%%%%


%%%%%%%%%%%%%%%%%%%%%%%%%%%%%%%%%%%%%%%%
\begin{quiz}{Conceptos de funciones vectoriales y curvas}
%%%%%%%%%%%%%%%%%%%%%%%%%%%%%%%%%%%%%%%%

\begin{multi}[]%
    {FV-Def-01}
    Sea $\func{f}{\Omega}{\R^m}$ con $\Omega\subset\R^n$. ¿Cuál de los siguientes casos corresponde a una {trayectoria}?
    \item* \(n=1\) y \(m>1\)
    \item \(n>1\) y $m=1$
    \item \(n>1\) y \(m>1\)
    \item $n=1$ y $m=1$
\end{multi}

\begin{multi}[]%
    {FV-Def-01b}
    La función $\func{f}{\R}{\R^3}$ definida por $f(t)=(t, t^2, \sen t)$ es:
    \item* Una trayectoria
    \item Un campo escalar
    \item Un campo vectorial
    \item Una función real de variable real
\end{multi}

\begin{multi}[]%
    {FV-Def-02}
    Sea $\func{f}{\Omega}{\R^m}$ con $\Omega\subset\R^n$. ¿Cuál de los siguientes casos corresponde a un {campo escalar}?
    \item* \(n>1\) y $m=1$
    \item $n=1$ y \(m>1\)
    \item \(n>1\) y \(m>1\)
    \item $n=1$ y $m=1$
\end{multi}

\begin{multi}[]%
    {FV-Def-02b}
    La función $\func{g}{\R^2}{\R}$ definida por $g(x,y)=x^2+y^2$ es:
    \item* Un campo escalar
    \item Una trayectoria
    \item Un campo vectorial
    \item Una función real de variable real
\end{multi}

\begin{multi}[]%
    {FV-NivelCurva-01}
    Dado un campo escalar $\func{f}{\Omega}{\R}$ con $\Omega\subset\R^2$ y $c\in\R$, ¿cómo se llama al conjunto $L_c(f)=\{x\in\Omega : f(x)=c\}$?
    \item* Curva de nivel
    \item Superficie de nivel
    \item Trayectoria de nivel
    \item Campo de nivel
\end{multi}

\begin{multi}[]%
    {FV-NivelCurva-01b}
    Para $\func{f}{\R^2}{\R}$ y $c\in\R$, el conjunto $\{(x,y)\in\R^2 : f(x,y)=3\}$ se denomina:
    \item* Curva de nivel de $f$ en $c=3$
    \item Superficie de nivel de $f$ en $c=3$
    \item Gráfica de $f$ en $c=3$
    \item Dominio de $f$
\end{multi}

\begin{multi}[]%
    {FV-VectorTangente-01}
    Dada una curva $C$ parametrizada por $\func{\alpha}{I}{\R^m}$, ¿cuál es el vector tangente a $C$ en el punto $\alpha(t)$, cuando existe y es no nulo?
    \item* $\alpha'(t)$
    \item $\alpha(t)$
    \item $\alpha''(t)$
    \item $\|\alpha'(t)\|$
\end{multi}

\begin{multi}[]%
    {FV-VectorTangente-01b}
    Si $\alpha'(t_0)\neq 0$, la recta tangente a la curva $C$ en el punto $\alpha(t_0)$ tiene como vector director:
    \item* $\alpha'(t_0)$
    \item $\alpha(t_0)$
    \item $\alpha''(t_0)$
    \item $\alpha(t_0)+\alpha'(t_0)$
\end{multi}

\begin{multi}[]%
    {FV-Movimiento-01}
    Si $\alpha$ describe el movimiento de una partícula, ¿cuál de las siguientes afirmaciones es correcta?
    \item* $\alpha'(t)$ es la velocidad y $\|\alpha'(t)\|$ es la rapidez en el tiempo $t$.
    \item $\alpha'(t)$ es la aceleración en el tiempo $t$.
    \item $\alpha''(t)$ es la velocidad en el tiempo $t$.
    \item $\|\alpha'(t)\|$ es la aceleración escalar en el tiempo $t$.
\end{multi}

\begin{multi}[]%
    {FV-Movimiento-01b}
    Si $\alpha$ describe la posición de una partícula, ¿cuál de las siguientes expresiones representa la aceleración en el tiempo $t$?
    \item* $\alpha''(t)$
    \item $\alpha'(t)$
    \item $\|\alpha'(t)\|$
    \item $\alpha(t)$
\end{multi}

\begin{multi}[]%
    {FV-CadenaRegla-01}
    Si $\func{u}{I}{\R}$ es derivable en $a$ y $\func{\alpha}{J}{\R^m}$ es derivable en $u(a)$, la regla de la cadena para $\alpha\circ u$ establece que:
    \item* $(\alpha\circ u)'(a)=u'(a)\,\alpha'(u(a))$
    \item $(\alpha\circ u)'(a)=\alpha'(a)\,u'(u(a))$
    \item $(\alpha\circ u)'(a)=\alpha'(u(a))+u'(a)$
    \item $(\alpha\circ u)'(a)=u'(a)+\alpha'(a)$
\end{multi}


\end{quiz}


%%%%%%%%%%%%%%%%%%%%%%%%%%%%%%%%%%%%%%%%
\begin{quiz}{Geometría de curvas}
%%%%%%%%%%%%%%%%%%%%%%%%%%%%%%%%%%%%%%%%

\begin{multi}[]%
    {GC-VectoresTriedro-01}
    Dada $\func{\alpha}{I}{\R^3}$ con $\alpha'(t)\neq 0$, ¿cuál es la definición del vector unitario tangente $T_\alpha(t)$?
    \item* $T_\alpha(t)=\dfrac{\alpha'(t)}{\|\alpha'(t)\|}$
    \item $T_\alpha(t)=\dfrac{T_\alpha'(t)}{\|T_\alpha'(t)\|}$
    \item $T_\alpha(t)=T_\alpha(t)\times N_\alpha(t)$
    \item $T_\alpha(t)=\alpha'(t)\cdot\alpha''(t)$
\end{multi}

\begin{multi}[]%
    {GC-VectoresTriedro-01b}
    El vector $T_\alpha(t)$ es unitario porque:
    \item* Se obtiene dividiendo $\alpha'(t)$ entre su norma $\|\alpha'(t)\|$.
    \item $\alpha'(t)$ siempre tiene norma 1 si la curva está parametrizada.
    \item Se define como el producto vectorial $T_\alpha(t)\times N_\alpha(t)$.
    \item Se obtiene integrando $\alpha''(t)$.
\end{multi}

\begin{multi}[]%
    {GC-VectoresTriedro-02}
    El vector binormal principal $B_\alpha(t)$ se define como:
    \item* $B_\alpha(t)=T_\alpha(t)\times N_\alpha(t)$
    \item $B_\alpha(t)=N_\alpha(t)\times T_\alpha(t)$
    \item $B_\alpha(t)=\dfrac{T_\alpha'(t)}{\|T_\alpha'(t)\|}$
    \item $B_\alpha(t)=\dfrac{\alpha'(t)}{\|\alpha'(t)\|}$
\end{multi}

\begin{multi}[]%
    {GC-VectoresTriedro-02b}
    El vector normal principal $N_\alpha(t)$ se define como:
    \item* $N_\alpha(t)=\dfrac{T_\alpha'(t)}{\|T_\alpha'(t)\|}$
    \item $N_\alpha(t)=T_\alpha(t)\times B_\alpha(t)$
    \item $N_\alpha(t)=\dfrac{\alpha'(t)}{\|\alpha'(t)\|}$
    \item $N_\alpha(t)=\alpha''(t)$
\end{multi}

\begin{multi}[]%
    {GC-Planos-01}
    ¿Cuál de los siguientes planos contiene a los vectores $T_\alpha(t)$ y $N_\alpha(t)$?
    \item* Plano osculador
    \item Plano normal
    \item Plano rectificante
    \item Plano tangente
\end{multi}

\begin{multi}[]%
    {GC-Planos-01b}
    ¿Cuál de los siguientes planos contiene a los vectores $B_\alpha(t)$ y $N_\alpha(t)$?
    \item* Plano normal
    \item Plano osculador
    \item Plano rectificante
    \item Plano tangente
\end{multi}

\begin{multi}[]%
    {GC-Curvatura-01}
    La curvatura $\kappa_\alpha(t)$ se define como:
    \item* $\kappa_\alpha(t)=\dfrac{\|T_\alpha'(t)\|}{\rho_\alpha(t)}$
    \item $\kappa_\alpha(t)=\|T_\alpha(t)\|$
    \item $\kappa_\alpha(t)=\rho_\alpha(t)\,\|T_\alpha'(t)\|$
    \item $\kappa_\alpha(t)=\dfrac{\rho_\alpha(t)}{\|T_\alpha'(t)\|}$
\end{multi}

\begin{multi}[]%
    {GC-Curvatura-01b}
    Si $\kappa_\alpha(t)=0$ para todo $t\in I$, ¿qué forma tiene la curva?
    \item* Es una recta (segmento de línea recta).
    \item Es una circunferencia.
    \item Es una hélice.
    \item Es una curva plana genérica.
\end{multi}

\begin{multi}[]%
    {GC-Curvatura-02}
    La curvatura también puede calcularse con la fórmula alternativa:
    \item* $\kappa_\alpha(t)=\dfrac{\|a_\alpha(t)\times v_\alpha(t)\|}{\rho^3_\alpha(t)}$
    \item $\kappa_\alpha(t)=\dfrac{\|v_\alpha(t)\times a_\alpha(t)\|}{\rho_\alpha(t)}$
    \item $\kappa_\alpha(t)=\dfrac{v_\alpha(t)\cdot a_\alpha(t)}{\rho^2_\alpha(t)}$
    \item $\kappa_\alpha(t)=\dfrac{\|a_\alpha(t)\|}{\|v_\alpha(t)\|}$
\end{multi}

\begin{multi}[]%
    {GC-Torsion-01}
    La torsión $\tau_\alpha(t)$ se define como:
    \item* $\tau_\alpha(t)=\dfrac{v_\alpha(t)\cdot\left(a_\alpha(t)\times a_\alpha'(t)\right)}{\|v_\alpha(t)\times a_\alpha(t)\|^2}$
    \item $\tau_\alpha(t)=\dfrac{\|T_\alpha'(t)\|}{\rho_\alpha(t)}$
    \item $\tau_\alpha(t)=\dfrac{a_\alpha(t)\cdot v_\alpha(t)}{\rho^3_\alpha(t)}$
    \item $\tau_\alpha(t)=\dfrac{\|v_\alpha(t)\times a_\alpha(t)\|}{v_\alpha(t)\cdot a_\alpha'(t)}$
\end{multi}

\begin{multi}[]%
    {GC-Torsion-01b}
    Una curva plana (contenida en un plano) tiene torsión:
    \item* $\tau=0$ en todos sus puntos.
    \item $\tau=1$ en todos sus puntos.
    \item $\tau$ igual a su curvatura en cada punto.
    \item $\tau$ indefinida, pues la torsión requiere que la curva esté en $\R^3$.
\end{multi}

\begin{multi}[]%
    {GC-LongitudArco-01}
    La longitud de arco de la curva descrita por $\func{\alpha}{[a,b]}{\R^n}$ está dada por:
    \item* $\ell(C)=\displaystyle\int_a^b \rho_\alpha(t)\,dt$
    \item $\ell(C)=\displaystyle\int_a^b \|\alpha(t)\|\,dt$
    \item $\ell(C)=\displaystyle\int_a^b \alpha'(t)\,dt$
    \item $\ell(C)=\displaystyle\int_a^b \kappa_\alpha(t)\,dt$
\end{multi}

\end{quiz}


%%%%%%%%%%%%%%%%%%%%%%%%%%%%%%%%%%%%%%%%
\begin{quiz}{Tipo de funciones}
%%%%%%%%%%%%%%%%%%%%%%%%%%%%%%%%%%%%%%%%

\begin{multi}[]%
    {TF-Clasif-01}
    La función $\func{f}{\R^2}{\R^3}$ es:
    \item* Un campo vectorial
    \item Una trayectoria
    \item Un campo escalar
\end{multi}

\begin{multi}[]%
    {TF-Clasif-01b}
    La función $\func{f}{\R}{\R^2}$ es:
    \item* Una trayectoria
    \item Un campo escalar
    \item Un campo vectorial
\end{multi}

\begin{multi}[]%
    {TF-Clasif-02}
    La función $\func{f}{\R^3}{\R}$ es:
    \item* Un campo escalar
    \item Un campo vectorial
    \item Una trayectoria
\end{multi}

\begin{multi}[]%
    {TF-Clasif-02b}
    La función $\func{F}{\R^2}{\R^2}$ definida por $F(x,y)=(x-y,\, x+y)$ es:
    \item* Un campo vectorial
    \item Un campo escalar
    \item Una trayectoria
\end{multi}

\begin{multi}[]%
    {TF-Clasif-03}
    La función definida por $f(t)=(t^2, t)$ es:
    \item* Una trayectoria
    \item Un campo escalar
    \item Un campo vectorial
\end{multi}

\begin{multi}[]%
    {TF-Clasif-03b}
    La función definida por $g(x,y)=\sqrt{x^2+y^2}$ es:
    \item* Un campo escalar
    \item Una trayectoria
    \item Un campo vectorial
\end{multi}

\begin{multi}[]%
    {TF-Clasif-04}
    La función definida por $f(t)=(t, \cos t, \sen t)$ es:
    \item* Una trayectoria
    \item Un campo escalar
    \item Un campo vectorial
\end{multi}

\begin{multi}[]%
    {TF-Clasif-04b}
    La función $\func{h}{\R^3}{\R^3}$ definida por $h(x,y,z)=(x+y,\, y+z,\, x+z)$ es:
    \item* Un campo vectorial
    \item Un campo escalar
    \item Una trayectoria
\end{multi}

\begin{multi}[]%
    {TF-Clasif-05}
    La función definida por $f(x,y)=x^2+y^2$ es:
    \item* Un campo escalar
    \item Una trayectoria
    \item Un campo vectorial
\end{multi}

\begin{multi}[]%
    {TF-Clasif-05b}
    La función $F(x,y,z)=x^2+y^2+z^2$ es:
    \item* Un campo escalar
    \item Una trayectoria
    \item Un campo vectorial
\end{multi}

\begin{multi}[]%
    {TF-Clasif-06}
    La función definida por $F(x,y)=(x+y,\, x-y,\, xy)$ es:
    \item* Un campo vectorial
    \item Una trayectoria
    \item Un campo escalar
\end{multi}

\begin{multi}[]%
    {TF-Clasif-06b}
    La función $\func{G}{\R^2}{\R^3}$ es:
    \item* Un campo vectorial
    \item Una trayectoria
    \item Un campo escalar
\end{multi}

\end{quiz}


%%%%%%%%%%%%%%%%%%%%%%%%%%%%%%%%%%%%%%%%
\begin{quiz}{Gráficas de funciones}
%%%%%%%%%%%%%%%%%%%%%%%%%%%%%%%%%%%%%%%%

\begin{multi}[]%
    {GF-Graf-01}
    La representación gráfica de la función $\alpha(t)=(t,\, t)$ es:
    \item* Una curva en el plano
    \item Una curva en el espacio
    \item Una superficie
    \item Un campo de flechas en el plano
\end{multi}

\begin{multi}[]%
    {GF-Graf-01b}
    La representación gráfica de la función $\alpha(t)=(t^2,\, t^3)$ es:
    \item* Una curva en el plano
    \item Una curva en el espacio
    \item Una superficie en $\R^3$
    \item Un campo de flechas en el plano
\end{multi}

\begin{multi}[]%
    {GF-Graf-02}
    La representación gráfica de la función $\alpha(t)=(t,\, t^2,\, t^3)$ es:
    \item* Una curva en el espacio
    \item Una curva en el plano
    \item Una superficie en $\R^3$
    \item Un campo de flechas en $\R^3$
\end{multi}

\begin{multi}[]%
    {GF-Graf-02b}
    La representación gráfica de la función $\alpha(t)=(\cos t,\, \sen t,\, t)$ es:
    \item* Una curva en el espacio (hélice)
    \item Una curva en el plano
    \item Una superficie en $\R^3$
    \item Un campo de flechas en $\R^3$
\end{multi}

\begin{multi}[]%
    {GF-Graf-03}
    La representación gráfica de la función $f(x,y)=x^2+y^2$ es:
    \item* Una superficie en $\R^3$
    \item Una curva en el plano
    \item Una curva en el espacio
    \item Un campo de flechas en $\R^2$
\end{multi}

\begin{multi}[]%
    {GF-Graf-03b}
    La representación gráfica de la función $f(x,y)=\sen(x)\cos(y)$ es:
    \item* Una superficie en $\R^3$
    \item Una curva en el plano
    \item Una curva en el espacio
    \item Un campo de flechas en $\R^2$
\end{multi}

\begin{multi}[]%
    {GF-Graf-04}
    La representación gráfica de la función $F(x,y)=(-y,\, x)$ es:
    \item* Un campo de flechas en el plano
    \item Una curva en el plano
    \item Una superficie en $\R^3$
    \item Una curva en el espacio
\end{multi}

\begin{multi}[]%
    {GF-Graf-04b}
    La representación gráfica de la función $F(x,y,z)=(x,\, y,\, z)$ es:
    \item* Un campo de flechas en el espacio
    \item Una curva en el espacio
    \item Una superficie en $\R^3$
    \item Una curva en el plano
\end{multi}

\begin{multi}[]%
    {GF-Graf-05}
    La representación gráfica de la función $\alpha(t)=(\cos t,\, \sen t)$ es:
    \item* Una curva en el plano (una circunferencia)
    \item Una curva en el espacio
    \item Una superficie en $\R^3$
    \item Un punto en el plano
\end{multi}

\begin{multi}[]%
    {GF-Graf-05b}
    La representación gráfica de la función $\alpha(t)=(\cos t,\, \sen t,\, 0)$, $t\in[0,2\pi]$, es:
    \item* Una curva en el espacio (circunferencia en el plano $xy$)
    \item Una curva en el plano
    \item Una superficie en $\R^3$
    \item Un campo de flechas en $\R^3$
\end{multi}

\end{quiz}


%%%%%%%%%%%%%%%%%%%%%%%%%%%%%%%%%%%%%%%%
\begin{quiz}{Cálculo en trayectorias}
%%%%%%%%%%%%%%%%%%%%%%%%%%%%%%%%%%%%%%%%

\begin{numerical}[tolerance=0.01]%
    {CT-Deriv-01}
    Sea $\alpha(t)=(t^2,\, t^3)$. La primera componente de $\alpha'(1)$ es:
    \item 2
\end{numerical}

\begin{numerical}[tolerance=0.01]%
    {CT-Deriv-01b}
    Sea $\alpha(t)=(3t,\, t^2)$. La segunda componente de $\alpha'(2)$ es:
    \item 4
\end{numerical}

\begin{numerical}[tolerance=0.01]%
    {CT-Deriv-02}
    Sea $\alpha(t)=(\sen t,\, \cos t)$. La segunda componente de $\alpha'(0)$ es:
    \item 0
\end{numerical}

\begin{numerical}[tolerance=0.01]%
    {CT-Deriv-02b}
    Sea $\alpha(t)=(e^t,\, \ln(t+1))$. La primera componente de $\alpha'(0)$ es:
    \item 1
\end{numerical}

\begin{numerical}[tolerance=0.01]%
    {CT-Deriv-03}
    Sea $\alpha(t)=(e^t,\, t^2,\, 3t)$. La tercera componente de $\alpha'(0)$ es:
    \item 3
\end{numerical}

\begin{numerical}[tolerance=0.01]%
    {CT-Deriv-03b}
    Sea $\alpha(t)=(\cos t,\, \sen t,\, t^2)$. La segunda componente de $\alpha'(0)$ es:
    \item 1
\end{numerical}

\begin{numerical}[tolerance=0.01]%
    {CT-Integ-01}
    Sea $\alpha(t)=(2t,\, 1)$. La primera componente de $\displaystyle\int_0^1 \alpha(t)\,dt$ es:
    \item 1
\end{numerical}

\begin{numerical}[tolerance=0.01]%
    {CT-Integ-01b}
    Sea $\alpha(t)=(t^2,\, 3)$. La segunda componente de $\displaystyle\int_0^2 \alpha(t)\,dt$ es:
    \item 6
\end{numerical}

\begin{numerical}[tolerance=0.01]%
    {CT-Integ-02}
    Sea $\alpha(t)=(t,\, t^2)$. La segunda componente de $\displaystyle\int_0^1 \alpha(t)\,dt$ es:
    \item 0.333
\end{numerical}

\begin{numerical}[tolerance=0.01]%
    {CT-Integ-02b}
    Sea $\alpha(t)=(t,\, e^t)$. La segunda componente de $\displaystyle\int_0^1 \alpha(t)\,dt$ es:
    \item 1.718
\end{numerical}

\end{quiz}


%%%%%%%%%%%%%%%%%%%%%%%%%%%%%%%%%%%%%%%%
\begin{quiz}{Tangentes a trayectorias}
%%%%%%%%%%%%%%%%%%%%%%%%%%%%%%%%%%%%%%%%

\begin{multi}[]%
    {TT-Tangente-01}
    El vector tangente a la trayectoria $\alpha(t)=(t^2,\, 2t)$ en $t=1$ es:
    \item* $(2,\, 2)$
    \item $(1,\, 2)$
    \item $(2,\, 1)$
    \item $(4,\, 2)$
\end{multi}

\begin{multi}[]%
    {TT-Tangente-01b}
    El vector tangente a la trayectoria $\alpha(t)=(t^3,\, t^2+1)$ en $t=1$ es:
    \item* $(3,\, 2)$
    \item $(1,\, 2)$
    \item $(3,\, 1)$
    \item $(1,\, 1)$
\end{multi}

\begin{multi}[]%
    {TT-Tangente-02}
    El vector tangente a la trayectoria $\alpha(t)=(\cos t,\, \sen t,\, t)$ en $t=0$ es:
    \item* $(0,\, 1,\, 1)$
    \item $(1,\, 0,\, 0)$
    \item $(-1,\, 0,\, 1)$
    \item $(0,\, 0,\, 1)$
\end{multi}

\begin{multi}[]%
    {TT-Tangente-02b}
    La dirección del vector tangente a la trayectoria $\alpha(t)=(t,\, t^2)$ en el punto $\alpha(1)=(1,1)$ es:
    \item* $(1,\, 2)$
    \item $(1,\, 1)$
    \item $(0,\, 2)$
    \item $(2,\, 1)$
\end{multi}

\begin{numerical}[tolerance=0.01]%
    {TT-TangenteNum-01}
    Sea $\alpha(t)=(t^3,\, t^2)$. La primera coordenada del vector tangente en $t=2$ es:
    \item 12
\end{numerical}

\begin{numerical}[tolerance=0.01]%
    {TT-TangenteNum-01b}
    Sea $\alpha(t)=(t^2,\, t^4)$. La segunda coordenada del vector tangente en $t=1$ es:
    \item 4
\end{numerical}

\begin{numerical}[tolerance=0.01]%
    {TT-TangenteNum-02}
    Sea $\alpha(t)=(\sen t,\, \cos t,\, 2t)$. La tercera coordenada del vector tangente en $t=0$ es:
    \item 2
\end{numerical}

\begin{numerical}[tolerance=0.01]%
    {TT-TangenteNum-02b}
    Sea $\alpha(t)=(t^2,\, e^t,\, \cos t)$. La primera coordenada del vector tangente en $t=0$ es:
    \item 0
\end{numerical}

\begin{multi}[]%
    {TT-Tangente-03}
    El vector tangente a la trayectoria $\alpha(t)=(e^t,\, e^{-t})$ en $t=0$ es:
    \item* $(1,\, -1)$
    \item $(1,\, 1)$
    \item $(0,\, 0)$
    \item $(-1,\, 1)$
\end{multi}

\begin{multi}[]%
    {TT-Tangente-03b}
    El vector tangente a la trayectoria $\alpha(t)=(\sen t,\, t,\, t^2)$ en $t=0$ es:
    \item* $(1,\, 1,\, 0)$
    \item $(0,\, 1,\, 0)$
    \item $(1,\, 0,\, 0)$
    \item $(0,\, 0,\, 0)$
\end{multi}

\end{quiz}


%%%%%%%%%%%%%%%%%%%%%%%%%%%%%%%%%%%%%%%%
\begin{quiz}{Ideas de curvatura y torsión}
%%%%%%%%%%%%%%%%%%%%%%%%%%%%%%%%%%%%%%%%

\begin{multi}[]%
    {ICT-Curv-01}
    Intuitivamente, ¿qué mide la curvatura $\kappa$ de una curva en un punto?
    \item* Qué tanto se dobla o se curva la trayectoria en ese punto.
    \item La velocidad con que se recorre la curva.
    \item La longitud total de la curva.
    \item El ángulo que forma la curva con el eje horizontal.
\end{multi}

\begin{multi}[]%
    {ICT-Curv-01b}
    ¿Qué ocurre con la curvatura $\kappa$ en un punto donde la trayectoria es casi recta?
    \item* La curvatura es pequeña, cercana a cero.
    \item La curvatura es grande, mayor que 1.
    \item La curvatura es igual a la rapidez en ese punto.
    \item La curvatura no está definida en ese punto.
\end{multi}

\begin{multi}[]%
    {ICT-Curv-02}
    Una línea recta tiene curvatura:
    \item* $\kappa=0$ en todos sus puntos.
    \item $\kappa=1$ en todos sus puntos.
    \item $\kappa$ variable según el punto.
    \item $\kappa$ indefinida, pues no se puede calcular en rectas.
\end{multi}

\begin{multi}[]%
    {ICT-Curv-03}
    Una circunferencia de radio $r$ grande comparada con una de radio $r$ pequeño tiene:
    \item* Menor curvatura, pues se curva menos.
    \item Mayor curvatura, pues es más larga.
    \item La misma curvatura, pues ambas son circunferencias.
    \item Curvatura cero, pues en el límite se parece a una recta.
\end{multi}

\begin{multi}[]%
    {ICT-Curv-03b}
    Si el radio de una circunferencia aumenta, su curvatura:
    \item* Disminuye, pues la curva se dobla menos.
    \item Aumenta, pues la circunferencia es más grande.
    \item Se mantiene constante, pues la curvatura no depende del radio.
    \item Se hace cero cuando el radio supera 1.
\end{multi}

\begin{multi}[]%
    {ICT-Tors-01}
    Intuitivamente, ¿qué mide la torsión $\tau$ de una curva en el espacio?
    \item* Qué tanto la curva se aparta del plano osculador, es decir, cuánto «se tuerce» fuera de ese plano.
    \item Qué tanto se dobla la curva en ese punto.
    \item La velocidad escalar de la partícula que recorre la curva.
    \item La longitud de arco acumulada hasta ese punto.
\end{multi}

\begin{multi}[]%
    {ICT-Tors-01b}
    Una curva tiene torsión nula si y solo si:
    \item* La curva está contenida en un plano.
    \item La curvatura es constante.
    \item La curva es una recta.
    \item La rapidez de recorrido es constante.
\end{multi}

\begin{multi}[]%
    {ICT-Tors-02}
    Una curva plana (contenida en un plano) tiene torsión:
    \item* $\tau=0$ en todos sus puntos.
    \item $\tau=1$ en todos sus puntos.
    \item $\tau$ variable que depende de la curvatura.
    \item $\tau$ indefinida, pues la torsión solo se define en el espacio.
\end{multi}

\begin{multi}[]%
    {ICT-Tors-02b}
    ¿Cuál de las siguientes afirmaciones sobre la torsión es correcta?
    \item* La torsión mide cuánto la curva se aparta del plano osculador en cada punto.
    \item La torsión mide cuánto se dobla la curva en cada punto.
    \item La torsión coincide siempre con la curvatura.
    \item La torsión solo está definida para curvas planas.
\end{multi}

\begin{multi}[]%
    {ICT-Tors-03}
    ¿Cuál de las siguientes curvas tiene torsión no nula?
    \item* Una hélice circular $\alpha(t)=(\cos t,\, \sen t,\, t)$.
    \item Una circunferencia en el plano $xy$.
    \item Una parábola en el plano $xz$.
    \item Una recta en $\R^3$.
\end{multi}

\begin{multi}[]%
    {ICT-Tors-03b}
    La hélice $\alpha(t)=(\cos t,\, \sen t,\, t)$ tiene curvatura y torsión:
    \item* Ambas constantes y no nulas.
    \item Curvatura constante y torsión nula.
    \item Curvatura nula y torsión constante.
    \item Ambas variables según el punto.
\end{multi}

\end{quiz}

\end{document}
